% Candidate 3: annotated heat map (a transfer matrix: trained on the row domain, tested on the column domain).
% Single-hue colour map from white to the paper's blue, values printed in every cell, white text on dark cells,
% the diagonal outlined, a slim colour bar.
\documentclass[10pt,tikz,border=2pt]{standalone}
\usepackage{plotfig}
\begin{document}
\begin{tikzpicture}
\begin{axis}[paper, width=4.6cm, height=4.6cm,
  axis line style={draw=none}, tick style={draw=none}, ymajorgrids=false,
  xtick=data, ytick=data, xticklabels={news, legal, bio, code, chat, math}, yticklabels={news, legal, bio, code, chat, math},
  x tick label style={text=ink, font=\fontsize{7}{8}\selectfont, yshift=1pt}, y tick label style={text=ink, font=\fontsize{7}{8}\selectfont},
  xlabel={tested on}, ylabel={trained on}, y dir=reverse, enlargelimits=false,
  colormap={paperblues}{color(0cm)=(white); color(1cm)=(kept)}, point meta min=0, point meta max=100,
  colorbar, colorbar style={width=0.16cm, ytick={0,25,50,75,100}, axis line style={draw=none}, tick style={draw=none},
    yticklabel={\pgfmathprintnumber{\tick}}, /pgf/number format/assume math mode=true,
    yticklabel style={font=\fontsize{6}{6}\selectfont, text=black!60, xshift=-1pt}},
  visualization depends on={value \thisrow{v} \as \cellv},   % "value" keeps the cell text as written; without it pgfmath prints 91.000000
  nodes near coords={\pgfmathtruncatemacro{\ci}{\cellv}\ifnum\ci>55 \color{white}\fi\cellv},
  every node near coord/.append style={anchor=center, font=\fontsize{6.2}{6}\selectfont\ttfamily, text=black!80}]
\addplot[matrix plot*, mesh/cols=6, point meta=explicit] table[meta=v] {
x y v
0 0 91
1 0 62
2 0 48
3 0 21
4 0 57
5 0 30
0 1 58
1 1 88
2 1 44
3 1 19
4 1 41
5 1 27
0 2 47
1 2 46
2 2 90
3 2 24
4 2 39
5 2 35
0 3 25
1 3 22
2 3 28
3 3 93
4 3 33
5 3 52
0 4 61
1 4 45
2 4 42
3 4 31
4 4 86
5 4 34
0 5 29
1 5 24
2 5 33
3 5 49
4 5 36
5 5 89
};
% paths inside an axis are drawn at \end{axis}, so a \foreach variable is gone by then; \pgfplotsinvokeforeach substitutes #1 now
\pgfplotsinvokeforeach{0.5,1.5,2.5,3.5,4.5}{
  \draw[white, line width=1pt] (axis cs:#1,-0.5) -- (axis cs:#1,5.5);
  \draw[white, line width=1pt] (axis cs:-0.5,#1) -- (axis cs:5.5,#1);}
\pgfplotsinvokeforeach{0,1,2,3,4,5}{
  \draw[ink, line width=0.6pt] (axis cs:#1-0.5,#1-0.5) rectangle (axis cs:#1+0.5,#1+0.5);}
\end{axis}
\end{tikzpicture}
\end{document}
